\section{Edge circulation and its physical completion}
\label{ct:section}

The absolute-bypass condition permits a qualitatively different mechanism
when $\delta>\eta$: positive circulation inside a sharp section edge may
be accompanied by an equal negative return immediately outside it. This
section proves an exact electrostatic cross-flux calculation and a
finite-dimensional Maxwell jet-interpolation result for that mechanism.
It then states the additional global estimate needed to make the
mechanism an admissible counterexample. Theorems~\ref{cl:remainder}
and~\ref{rg:physical-limit} discharge that estimate in the following
sections, leading to Theorem~\ref{uc:counterexample}.
In particular, no impressed dipole is included in the admissible material
class and no electrostatic field is presented as an exact solution at the
fixed nonzero frequency.

The proposed mechanism exploits the difference between signed passage and
absolute bypass. A small inclusion close to the edge creates a local pattern
with forward flux on one side and return flux on the other. By assigning
independent local field data to different edge points, one can make this
pattern for every coherent input in a finite-dimensional space. That is a
local calculation: the device must still generate those data and keep the
error small on the entire plane. The next two sections supply those global
ingredients for the particular family used in the final theorem.

\subsection{Exact local profile and physical coherent jets}

Use the physical fields, the convention $e^{-i\omega t}$, and
$Z_0=\sqrt{\mu_0/\eps_0}$. At one straight edge introduce orthonormal
coordinates $(x,y,z)$ with the section locally $x<0$, outward normal
$\widehat x$, tangent $\widehat y$, and
$\widehat x\times\widehat y=\widehat z$. Fix dimensionless numbers
$0<\rho<h$ and $1<\eps_p\leq\eps_+$. Set
\begin{equation}\label{ct:alpha}
 \alpha=4\pi\rho^3\frac{\eps_p-1}{\eps_p+2}>0.
\end{equation}

\begin{lemma}[PROVED: electrostatic edge cross-flux]
\label{ct:profile}
For $s>0$, put a dielectric sphere of radius $\rho s$ at
$(0,0,hs)$. Its exact electrostatic exterior response to a uniform
normal electric field $A e\widehat z$, where $A>0$ and $e\in\mathbb C$,
has dipole moment $\eps_0\alpha s^3 A e\widehat z$.
On the plane $z=0$ its scattered tangential electric field is
\begin{equation}\label{ct:dipoletrace}
 E^{\rm stat}_t(s\xi,s\upsilon,0)
 =-\frac{3\alpha h A e}{4\pi}
       \frac{(\xi,\upsilon)}{(\xi^2+\upsilon^2+h^2)^{5/2}}.
\end{equation}
Cross this field with a constant tangential magnetic jet
$A b\widehat y$, $b\in\mathbb C$. The resulting real cross-flux
profile is exactly
\begin{equation}\label{ct:profileformula}
 f_s(s\xi,s\upsilon)
 :=\tfrac12\operatorname{Re}
    (E^{\rm stat}_t\times A\overline b\widehat y)\cdot\widehat z
 =-\frac{3\alpha h A^2}{8\pi}\operatorname{Re}(e\overline b)
       \frac{\xi}{(\xi^2+\upsilon^2+h^2)^{5/2}}.
\end{equation}
If $q=\operatorname{Re}(e\overline b)\geq0$, then
\begin{equation}\label{ct:profileintegrals}
 \int_{x<0}f_s\,dx\,dy
 =\int_{x>0}|f_s|\,dx\,dy
 =\frac{s^2 A^2\alpha q}{4\pi h},\qquad
 \int_{\mathbb R^2}f_s\,dx\,dy=0.
\end{equation}
The profile is nonnegative inside the edge and nonpositive outside.
These are exact statements about the stated electrostatic cross-term.
\end{lemma}
\begin{proof}
Center coordinates at the sphere. Write the exterior potential as
$-Ae z+cAe(\rho s)^3z/r^3$ and the interior potential as $-dAe z$.
Continuity of potential and normal electric displacement gives
$d=1-c$ and $\eps_p(1-c)=1+2c$. Therefore
$c=(\eps_p-1)/(\eps_p+2)$, proving the polarizability.
The dipole field is
$[3(p\cdot\widehat r)\widehat r-p]/(4\pi\eps_0 r^3)$.
At the observation plane the center-relative vector is
$s(\xi,\upsilon,-h)$, which gives \eqref{ct:dipoletrace}.
Taking its cross product gives \eqref{ct:profileformula}.
Finally,
\[
 \int_{\mathbb R}\frac{d\upsilon}
 { (\xi^2+\upsilon^2+h^2)^{5/2}}
 =\frac{4}{3(\xi^2+h^2)^2},\qquad
 \int_{-\infty}^0\frac{-\xi\,d\xi}{(\xi^2+h^2)^2}
 =\frac1{2h^2}.
\]
Together with $dx\,dy=s^2d\xi\,d\upsilon$ these prove the
half-plane integrals. Oddness in $\xi$ proves the last identity.
\end{proof}

An electric-wall vacuum Maxwell field satisfies
\eqref{cd:electric-wall}; in particular $E_t=H_z=0$ on the plane.
Its normal Poynting density is identically zero there, including for
every coherent combination of fields in that sector. We next prescribe
the nonzero jets $E_z,H_x,H_y$ while retaining this exact cancellation.

\begin{lemma}[PROVED: finite even-TE jet interpolation]
\label{ct:jets}
Let $p_1,\ldots,p_N$ be nonzero points of $z=0$, with
$p_i\ne p_j$ and $p_i\ne-p_j$ for $i\ne j$. Arbitrary complex values
of $(E_z,H_x,H_y)$ at these $N$ points can be realized by a finite
linear combination of regular vacuum TE spherical waves with even
angular order $l\geq2$ and even azimuthal index $m$.
The field is electric-wall and invariant, with the vectors rotated,
under a rotation by $\pi$ about the $z$ axis. In particular
\begin{equation}\label{ct:partnerjets}
 E_z(-p_j)=E_z(p_j),\qquad H_t(-p_j)=-H_t(p_j).
\end{equation}
The interpolating regular waves have finite physical incoming power.
No bound uniform in the interpolation points or target jets is asserted.
\end{lemma}
\begin{proof}
Consider a scalar regular Helmholtz solution $u$ that is even in $z$
and invariant under $(x,y)\mapsto(-x,-y)$, and set
\begin{equation}\label{ct:debye}
 E=\nabla\times(xu)=\nabla u\times x,
 \qquad H=\frac{\nabla\times E}{i\omega\mu_0}.
\end{equation}
Since $\nabla\cdot E=0$ and $(\Delta+k_0^2)E=0$, these fields
satisfy both vacuum Maxwell equations. They have precisely the
electric-wall and rotation symmetries just stated. Their scalar
spherical-harmonic expansion contains only even $m$ and even $l$;
the $l=0$ term gives zero in \eqref{ct:debye}. Thus finite truncations
are regular TE fields of the required kind.

The full tangential scalar two-jets can first be prescribed at
$p_1,\ldots,p_N$. To see this without assuming a Runge theorem,
write $r_t=(x,y)$ and use the entire scalar solutions
\[
 u_\xi(r_t,z)=\cos(\xi\cdot r_t)
            \cos\bigl(\sqrt{k_0^2-|\xi|^2}\,z\bigr),
 \qquad |\xi|<k_0.
\]
If a linear functional on those finitely many two-jets annihilates
every $u_\xi$, its value is an exponential polynomial in $\xi$,
vanishing on the real open disk. Its entire continuation therefore
vanishes identically. It is the Fourier transform of the symmetrized
finite distribution of derivatives of point masses supported at
$p_j$ and $-p_j$. These support points are all distinct, so uniqueness
of the Fourier transform forces every coefficient of the original
jet functional to be zero. Finite-dimensional duality proves the
surjectivity of the scalar jet map.

For clarity, the resulting scalar jets really do prescribe the three
Maxwell jets independently. Rotate coordinates in the plane so that
one point is $(r,0,0)$ with $r>0$. Formula \eqref{ct:debye} gives
\begin{equation}\label{ct:jetmap}
 E_z=-r u_y,\qquad
 H_x=\frac{2u_x+r u_{xx}+k_0^2r u}{i\omega\mu_0},\qquad
 H_y=\frac{2u_y+r u_{xy}}{i\omega\mu_0}.
\end{equation}
Choose $u_y$ for $E_z$, then $u_{xy}$ for $H_y$, and $u_{xx}$
for $H_x$. This gives a surjective map at each point, hence at all
points simultaneously.

Each $u_\xi$ is a finite sum of plane waves with real wavevectors
of length $k_0$. Replace its point-supported angular measures by
smooth, symmetrized densities on the sphere; their fields and any
fixed finite set of derivatives converge on compact sets. Thus the
jet map of smooth symmetric scalar Herglotz fields is onto: choose
finitely many approximately interpolating columns close enough to
an invertible interpolation matrix. For each of these smooth fields,
truncate the spherical-harmonic expansion, preserving both
symmetries. Compact derivative convergence again preserves the
invertibility of that finite matrix. Invert it to obtain exact jets
using one finite angular band. Its Maxwell fields are finite regular
TE sums; their incoming spherical amplitudes have finite power.
The parity identities give \eqref{ct:partnerjets}.
\end{proof}

\begin{lemma}[PROVED: balancing the physical input Gram matrix]
\label{ct:gram}
For any finite collection of target jets covered by Lemma~\ref{ct:jets},
there are interpolants $v_1,\ldots,v_N$ in a common finite even-TE
band whose physical incoming Gram matrix is $\lambda I_N$ for some
$0<\lambda<\infty$. Consequently $g_j=v_j/\sqrt\lambda$ are
orthonormal physical incoming data. All $v_j$ retain their prescribed
jets.
\end{lemma}
\begin{proof}
Start with arbitrary interpolants $u_j$ and their positive
semidefinite incoming Gram matrix $G$. The space of finite even-TE
regular waves has unbounded dimension, whereas the jet map has
finite rank. Enlarge its finite band until its jet kernel contains
an $N$-dimensional subspace orthogonal, in physical incoming power,
to the span of the $u_j$. Choose $\lambda>\|G\|$.
In an orthonormal basis of that kernel subspace, a matrix square
root of $\lambda I-G$ supplies vectors $w_j$ with Gram matrix
$\lambda I-G$. They have zero prescribed jets and are orthogonal
to all $u_j$. Hence $v_j=u_j+w_j$ has the same jets as $u_j$ and
Gram matrix $G+(\lambda I-G)=\lambda I$.
\end{proof}

\subsection{Finite-rank diagonal edge circulation}

Choose distinct points $p_j=(L/2,y_j,0)$ on the interior of the
right edge of the square, with $-L/2<y_j<L/2$. Their partners
$-p_j$ lie on the left edge. Let $\ell_0=k_0^{-1}$.
By Lemmas~\ref{ct:jets}--\ref{ct:gram}, take interpolants with jets
\begin{equation}\label{ct:unitjets}
 E_z^{v_j}(p_i)=\frac{\sqrt{Z_0}}{\ell_0}\delta_{ij},\qquad
 H_t^{v_j}(p_i)=\frac{\widehat y}{\sqrt{Z_0}\ell_0}\delta_{ij},
 \qquad \langle v_i,v_j\rangle_{\rm in}=\lambda\delta_{ij}.
\end{equation}
Physical field units are implicit in the unit-power normalization
of this display. Its relevant product is exactly $1/\ell_0^2$.

For each fixed amplification $A>0$, Theorem~\ref{cd:wall-isolation}
can amplify every even TE order in this finite band by the same
factor $A\sqrt\lambda$, up to a common phase and arbitrarily small
compact derivative errors. Applied to the orthonormal inputs
$g_j=v_j/\sqrt\lambda$, its ideal core jets at $p_j$ are
\begin{equation}\label{ct:amplifiedjets}
 E_z(c;p_j)=\frac{A\sqrt{Z_0}}{\ell_0}c_j,\qquad
 H_t(c;p_j)=\frac{A}{\sqrt{Z_0}\ell_0}c_j\widehat y,
 \qquad c\in\mathbb C^N.
\end{equation}
The original radial-device background has zero plane flux exactly,
at every finite stage, by reflection symmetry. Compact derivative
approximation concerns its illumination only, before adding spheres.

For $a>0$ set $s=a\ell_0$ and $A=t/a$, with fixed $t>0$.
At each of $p_j$ and $-p_j$, put the electrostatic cross-profile of
Lemma~\ref{ct:profile}, with the local outward/tangent directions
$(\widehat x,\widehat y)$ on the right and
$(-\widehat x,-\widehat y)$ on the left. Use the jets
\eqref{ct:amplifiedjets}, and denote the sum of these $2N$ profiles
by $F_a(c;x,y)$. Thus $F_a$ is defined on the whole plane by
explicit functions; it is quadratic in $c$ and diagonal in the
coefficients $|c_j|^2$.

\begin{proposition}[PROVED: coherent local-model passage and absolute bypass]
\label{ct:model}
For each fixed finite $N$, the profiles just defined obey
\begin{equation}\label{ct:modellimits}
 \sup_{\|c\|=1}\left|\int_D F_a(c)\,dA-p_0\right|\longrightarrow0,
 \qquad
 \sup_{\|c\|=1}\left|\int_{\mathbb R^2\setminus D}|F_a(c)|\,dA-p_0\right|
       \longrightarrow0,
 \qquad p_0=\frac{\alpha t^2}{2\pi h}.
\end{equation}
Also $\int_{\mathbb R^2}F_a(c)\,dA=0$ exactly. The geometric spheres
associated with these profiles have radius $\rho a\ell_0$, lie
strictly above the plane, are disjoint for sufficiently small $a$,
and preserve the rotation by $\pi$ symmetry. Their contrasts remain
in the fixed ordinary dielectric interval.
\end{proposition}
\begin{proof}
For one point, the product in Lemma~\ref{ct:profile} is
$q=|c_j|^2/\ell_0^2$, so its positive half-plane mass equals
$\alpha t^2|c_j|^2/(4\pi h)$. Its partner has the same mass:
both the local tangent and $H_t$ reverse at $-p_j$, leaving the
scalar product and the orientation of positive interior flux unchanged.
Summing gives $p_0\sum_j|c_j|^2$.

Choose disjoint fixed disks about the $2N$ edge points, avoiding
corners. Within each disk the square is exactly its interior
half-plane. The rescaled profile has an integrable envelope bounded
by a constant times $(1+\xi^2+\upsilon^2)^{-2}$, so its $L^1$ mass
outside its disk tends to zero. More precisely its tail outside
scaled radius $R$ is $O(R^{-2})$. On each disk retain only its own
profile. The $L^1$ difference between this truncated sum and $F_a$
tends to zero, uniformly for $\sum_j|c_j|^2=1$, since there are
finitely many centers. The retained profiles have the required
signs on the two sides. The inequality
$\bigl||u|-|v|\bigr|\leq|u-v|$ proves both limits, including the
absolute-value limit for arbitrary coherent $c$.
Oddness gives zero total signed integral for every individual
profile. Finally $h>\rho$ leaves a vacuum gap above the plane,
and the fixed distinct centers ensure disjointness when $a$ is small.
\end{proof}

\begin{remark}[Why conditioning of the interpolation is not a bypass loss]
\label{ct:conditioning}
The number $\lambda$ in Lemma~\ref{ct:gram} can be large and depends
on $N$ and the points. It is fixed before $a\downarrow0$ and is
absorbed into the common core gain $A\sqrt\lambda$. The incoming
data remain orthonormal. Thus the model's passage and bypass
matrices are both $p_0 I_N$ in physical input coordinates, rather
than matrices whose ratio deteriorates with an interpolation
condition number. This is a finite-dimensional statement and does
not supply a bound on the perturbed Maxwell Green operator.
\end{remark}

\subsection{The precise physical estimate still required}

For each $a$, the finite-band control theorem permits choosing a
finite radial background with vacuum core large enough that its
unperturbed fields approximate \eqref{ct:amplifiedjets} and any
prescribed finite number of derivatives as accurately as desired
on a fixed neighborhood of all the edge points. The target gain
$t\sqrt\lambda/a$ is finite for each $a$. This diagonal choice is
legitimate for illumination. Add the $2N$ actual dielectric spheres
to that background, and let $S_{z,a}(c)$ denote the plane flux of
the resulting exact total scattering field for the same orthonormal
incoming inputs. A background approximation theorem alone does
not estimate this new field.

\begin{proposition}[CONDITIONAL: a sufficient global completion estimate]
\label{ct:completion}
Fix $0<\tau<1$ and $0<\eta<\delta$. Choose $t$ so that the constant
$p_0$ of \eqref{ct:modellimits} lies strictly between $\eta$ and
$\delta$. Suppose that for each fixed $N$ one can choose the finite
radial backgrounds and sphere sizes above so that the actual
physical scattering fields have legitimate integrable plane fluxes
and satisfy
\begin{equation}\label{ct:missingestimate}
 \sup_{\|c\|=1}
 \|S_{z,a}(c)-F_a(c)\|_{L^1(\mathbb R^2)}\longrightarrow0.
\end{equation}
Take the selected output map to be the full unitary scattering map
of each finite lossless dielectric. Then for sufficiently small $a$
the span of those $N$ physical incoming data is counted under
Definition~\ref{pr:count}, with the same fixed $\tau,\eta,\delta$.
Consequently \emph{if} \eqref{ct:missingestimate} can be established
for every finite $N$, the unrestricted finiteness target is false
for this one admissible fixed threshold tuple.
\end{proposition}
\begin{proof}
Orthonormality gives input power $\|c\|^2$; unitarity gives
$\|Tc\|^2=\|c\|^2\geq\tau\|c\|^2$.
Equations \eqref{ct:modellimits}--\eqref{ct:missingestimate} imply,
uniformly for unit $c$, that actual passage tends to $p_0$ and
actual absolute bypass tends to $p_0$. Their strict margins from
$\eta$ and $\delta$ prove the other two counting inequalities.
Homogeneity covers all $c$. For each member the sphere-to-plane
gap is positive, so it supplies a source-free vacuum neighborhood
of the section; its shrinking size is not a forbidden resource.
\end{proof}

\begin{remark}[The bridge discharged in the following sections]
\label{ct:gap}
The exact isolated-sphere electrostatic polarizability and finite-band
illumination control alone do not imply \eqref{ct:missingestimate}.
Three issues must be controlled together:
the small-inclusion expansion in a background whose resonant
amplification grows like $a^{-1}$, the source Green response of
unselected angular orders, and the absolute flux of all correction
fields across the entire plane, including the remote layers.
A finite inverse for each separate device would be insufficient unless
its dependence on that varying background is controlled.

The sphere volume is $O(a^3)$ and the amplified jet is $O(a^{-1})$.
Theorem~\ref{cl:remainder} turns the resulting scale
$a^3a^{-2}=a$ into an actual bounded-inverse and whole-plane remainder
estimate for an explicit parity reaction kernel.
Theorem~\ref{rg:physical-limit} identifies that kernel as the source
response of distant finite common stacks in the limit, controlling
all angular orders and the remote absolute-flux tails. Its inverse
stability step transfers the estimate to actual perturbed devices.
These two additional proofs, rather than incident gain or signed
power conservation alone, discharge \eqref{ct:missingestimate}.
The final fixed-threshold construction is Theorem~\ref{uc:counterexample}.
\end{remark}
